% Zdroj: PDF priklady-leto-2023.pdf, fyzická str. 2. Značka: otázka studijního zaměření (neznačená starší sada). Zařazeno: S3.
\section{Testování binárního klasifikátoru}
\szzotazka{léto 2023}{6}{Testování binárního klasifikátoru}

\noindent\textbf{Zadání.}\par
Předpokládejme, že k testování binárního klasifikátoru byla použita testovací sada ve složení
$100$ pozitivních a $400$ negativních příkladů.

\begin{enumerate}
  \item[a)] Ukázalo se, že ROC křivka prochází bodem $\mathrm{TPR} = \mathrm{FPR} = 20\,\%$.
    Vypočítejte hodnotu precision v tomto bodě.
  \item[b)] Jiný klasifikátor vykazuje konstantní hodnotu $\text{precision} = 50\,\%$.
    Znázorněte odpovídající ROC křivku.
\end{enumerate}

\noindent Použité zkratky: ROC — Receiver Operating Characteristic curve; TPR — true-positive rate
(též: sensitivity neboli recall); FPR — false-positive rate (též: fall-out).

\medskip
\noindent\textbf{Řešení.} \textit{(V~PDF bez náčrtu řešení; vlastní vzorové řešení, výpočet ověřen numericky.)}\par
\begin{enumerate}
  \item[a)] $\mathrm{TPR}=0{,}2\Rightarrow \mathrm{TP}=0{,}2\cdot100=20$; $\mathrm{FPR}=0{,}2\Rightarrow \mathrm{FP}=0{,}2\cdot400=80$. Precision $=\dfrac{\mathrm{TP}}{\mathrm{TP}+\mathrm{FP}}=\dfrac{20}{100}=0{,}2$.
  \item[b)] Konstantní precision $\dfrac{\mathrm{TP}}{\mathrm{TP}+\mathrm{FP}}=\tfrac12$ znamená $\mathrm{FP}=\mathrm{TP}$ při každém prahu. Protože $\mathrm{TPR}=\mathrm{TP}/100$ a~$\mathrm{FPR}=\mathrm{FP}/400=\mathrm{TP}/400$, platí $\mathrm{TPR}=4\,\mathrm{FPR}$. ROC křivka je tedy \emph{úsečka} z~$(0,0)$ do bodu $(\mathrm{FPR},\mathrm{TPR})=(0{,}25,\,1)$ (leží nad diagonálou náhodného hádání).
\end{enumerate}
